\chapter{Vestibular schwannoma}
\index{Vestibular schwannoma}

\section*{Definition and Overview}
A vestibular schwannoma, historically and commonly referred to as an acoustic neuroma, is a benign, slow-growing, primary intracranial tumour that originates from the Schwann cells of the vestibular division of the vestibulocochlear nerve (cranial nerve VIII). The vestibulocochlear nerve is responsible for transmitting auditory and vestibular (balance) information from the inner ear to the brainstem. In rare instances, these neoplasms may arise from the cochlear division of the nerve, but they are overwhelmingly vestibular in origin. These tumours represent a significant clinical entity in neuro-otology and neurosurgery due to their anatomical position within the internal auditory canal (IAC) and their potential expansion into the cerebellopontine angle (CPA). The cerebellopontine angle is a critical fluid-filled space in the posterior cranial fossa bounded by the cerebellum, the pons, and the petrous temporal bone. It houses not only cranial nerve VIII but also cranial nerve VII (the facial nerve), cranial nerve V (the trigeminal nerve), and major regional vasculature, including the anterior inferior cerebellar artery (AICA).

While vestibular schwannomas are histologically benign (World Health Organization Grade I), meaning they lack metastatic potential and grow slowly, their location makes them clinically formidable. As the tumour expands, it compresses the surrounding neurovascular structures, leading to progressive hearing loss, tinnitus, vestibular dysfunction (such as dysequilibrium or vertigo), and facial nerve dysfunction. If left untreated, large tumours can exert significant mass effect on the brainstem and the cerebellum, leading to obstructive hydrocephalus, cerebellar ataxia, and potentially life-threatening elevations in intracranial pressure. Understanding the pathophysiology, clinical presentation, and diverse management strategies is critical for preventing irreversible neurological deficits and optimizing patient outcomes.

\section*{Epidemiology}
Vestibular schwannomas account for approximately 8\% of all primary intracranial neoplasms and nearly 80\% to 85\% of tumours located in the cerebellopontine angle. The annual incidence of clinically diagnosed vestibular schwannoma is estimated to be between 1.0 and 2.0 cases per 100,000 population in developed countries. However, with the widespread availability and refinement of high-resolution magnetic resonance imaging (MRI), the detection of small, asymptomatic, or minimally symptoms-producing incidental tumours has increased significantly over the past few decades, suggesting a higher subclinical prevalence.

The disease most commonly presents in individuals between the fourth and sixth decades of life, with a peak incidence around 50 to 60 years of age. There is generally no significant gender predilection, although some registry studies report a slight female predominance. Vestibular schwannomas are overwhelmingly unilateral and sporadic, representing approximately 95\% of all cases. Bilateral vestibular schwannomas are highly pathognomonic of Neurofibromatosis Type 2 (NF2), a rare autosomal dominant genetic disorder caused by mutations in the NF2 tumour suppressor gene on chromosome 22q12. Bilateral presentations typically occur in younger patients, often in childhood or early adulthood, and are associated with a more aggressive disease course, multiple other cranial and spinal tumours (such as meningiomas and ependymomas), and a poorer long-term prognosis.

\section*{Aetiology and Pathophysiology}
The underlying mechanisms driving the development of vestibular schwannomas are centered on the dysfunction of Schwann cells, which are responsible for producing the myelin sheath that insulates peripheral nerves, including the vestibular portion of cranial nerve VIII.
\begin{itemize}
    \item \textbf{Loss of the NF2 Gene Function:} The primary genetic alteration responsible for both sporadic and syndromic (NF2-associated) vestibular schwannomas is the inactivation of both alleles of the \textit{NF2} gene located on the long arm of chromosome 22 (22q12.2). In sporadic cases, this inactivation occurs due to somatic mutations or chromosome loss in both alleles within the Schwann cell line (a two-hit hypothesis model). In syndromic cases, one mutated allele is inherited germline, and the second allele is inactivated somatically.
    \item \textbf{Loss of Merlin (Schwannomin) Protein:} The \textit{NF2} gene encodes a membrane-cytoskeleton scaffolding protein called merlin (or schwannomin). Merlin acts as a tumour suppressor by facilitating contact-dependent inhibition of growth. When functional, merlin links actin filaments to cell membrane glycoproteins and suppresses mitogenic signaling pathways, particularly the Ras-MAPK, PI3K-Akt, and Hippo pathways. The loss of merlin results in uncontrolled Schwann cell proliferation, leading to tumorigenesis.
    \item \textbf{Anatomical Site of Origin:} The majority of vestibular schwannomas arise within the internal auditory canal, specifically at the Obersteiner-Redlich zone. This zone represents the transition point where the central myelin (produced by oligodendrocytes) transitions to peripheral myelin (produced by Schwann cells). This transition zone is highly susceptible to neoplastic transformation.
    \item \textbf{Macroscopic and Microscopic Growth Patterns:} Histologically, vestibular schwannomas exhibit a classic biphasic pattern consisting of Antoni A and Antoni B areas. \textbf{Antoni A areas} are highly cellular, containing spindle-shaped Schwann cells with elongated nuclei arranged in palisades or whorls around acellular eosinophilic regions (known as Verocay bodies). \textbf{Antoni B areas} are hypocellular, showing loose, myxoid tissue with pleomorphic cells and microcystic changes. The vascularity is often characterized by thick-walled, hyalinized blood vessels that are prone to thrombosis and microhaemorrhage.
    \item \textbf{Environmental and Occupational Risk Factors:} Extensive epidemiological investigations have failed to identify strong, consistent environmental triggers for sporadic vestibular schwannomas. Although prolonged exposure to high levels of loud noise (such as industrial noise) and cell phone radiation have been proposed as potential risk factors, the scientific consensus remains inconclusive. Therapeutic ionizing radiation to the head and neck region during childhood is an established, albeit rare, risk factor.
\end{itemize}

\section*{Clinical Features}
The clinical presentation of vestibular schwannomas is dictated by their size, location, and the specific cranial nerves and brain structures compressed by the expanding mass. The onset of symptoms is usually insidious, progressing over several years, which frequently delays diagnosis.
\begin{itemize}
    \item \textbf{Hearing Loss:} Unilateral or asymmetric progressive sensorineural hearing loss is the most common presenting symptom, occurring in over 90\% of patients. The hearing loss typically affects high frequencies first, impairing speech discrimination out of proportion to the pure-tone threshold loss. In approximately 5\% to 15\% of cases, patients may experience sudden sensorineural hearing loss, which is believed to result from labyrinthine artery occlusion, compression of the cochlear nerve, or sudden intratumoural haemorrhage.
    \item \textbf{Tinnitus:} Subjective, high-pitched unilateral tinnitus is present in approximately 70\% of patients. It often localizes to the affected ear and can precede the awareness of hearing loss by several years.
    \item \textbf{Vestibular Symptoms:} Because these tumours originate from the vestibular nerve, balance disturbances are common. However, due to the slow-growing nature of the tumour, the central vestibular system is often able to compensate, presenting as mild, chronic dysequilibrium, unsteadiness, or motion intolerance rather than acute, rotatory vertigo. True vertigo is less common, occurring in only 20\% to 30\% of cases, and is more frequently associated with smaller, rapidly growing, or labyrinthine-bound tumours.
    \item \textbf{Trigeminal Nerve Dysfunction:} As the tumour expands out of the internal auditory canal into the cerebellopontine angle, it can compress the trigeminal nerve (cranial nerve V). This leads to facial numbness, paraesthesia, hypoesthesia, and the loss of the corneal reflex. Corneal hypoesthesia is a classic diagnostic sign of a larger CPA mass.
    \item \textbf{Facial Nerve Dysfunction:} Despite the close anatomical proximity of the facial nerve (cranial nerve VII) to the vestibulocochlear nerve within the internal auditory canal, clinically apparent facial weakness is uncommon at initial presentation (occurring in less than 10\% of cases) because the motor fibres of the facial nerve are highly resistant to slow, progressive compression. When present, it may manifest as mild facial twitching (hemifacial spasm) or subtle weakness of the facial expression muscles.
    \item \textbf{Neurological Signs of Large Tumours:} When the tumour exceeds 3 cm to 4 cm in diameter, it exerts significant mass effect on the brainstem and cerebellum. This compression manifests as cerebellar ataxia (clumsiness, wide-based gait, intention tremor, nystagmus), dysphagia, hoarseness, and dysarthria due to compression of lower cranial nerves (IX, X, and XI). Ultimately, compression of the fourth ventricle impairs cerebrospinal fluid (CSF) flow, leading to obstructive hydrocephalus. This manifests as symptoms of raised intracranial pressure, including severe headache (often worse in the morning), nausea, projectile vomiting, papilloedema, and visual disturbances (such as diplopia due to abducens nerve palsy).
\end{itemize}

\section*{Differential Diagnosis}
When evaluating a patient with a cerebellopontine angle mass or progressive unilateral vestibulocochlear dysfunction, several other lesions must be considered.
\begin{itemize}
    \item \textbf{Meningioma:} Cerebellopontine angle meningiomas are the second most common CPA tumours, accounting for approximately 10\% to 15\% of cases. Unlike vestibular schwannomas, meningiomas typically arise from the dura mater of the posterior petrous bone, do not widen the internal auditory canal, and show a characteristic dural tail sign on contrast-enhanced MRI. They are also often eccentric to the IAC.
    \item \textbf{Epidermoid Cyst:} Epidermoid cysts (or primary cholesteatomas of the CPA) represent approximately 5\% of CPA lesions. They are congenital, ectodermal-derived slow-growing lesions that fill the CPA cisterns and envelop surrounding neurovascular structures rather than displacing them. On MRI, they are characteristically hyperintense on diffusion-weighted imaging (DWI) and hypointense on T1-weighted images without enhancement.
    \item \textbf{Arachnoid Cyst:} These are benign, fluid-filled sacs that contain cerebrospinal fluid. On MRI, they follow CSF signal intensity on all sequences, including FLAIR and DWI, and show no solid components or contrast enhancement.
    \item \textbf{Facial Schwannoma:} Schwannomas arising from the facial nerve (cranial nerve VII) are rare and can mimic vestibular schwannomas. They are clinically suggested when facial weakness is an early and prominent feature, or if the tumour extends along the anatomical course of the facial nerve (e.g., into the fallopian canal or geniculate ganglion).
    \item \textbf{Non-neoplastic Otological Diseases:} Diseases such as M\'{e}ni\`{e}re's disease, labyrinthitis, otosclerosis, and sudden sensorineural hearing loss due to viral or vascular etiologies can mimic the early audiometric and vestibular presentation of a small vestibular schwannoma.
    \item \textbf{Vascular Anomalies:} Ectatic loop of the anterior inferior cerebellar artery (AICA) or aneurysms within the cerebellopontine angle can compress the vestibulocochlear nerve, mimicking the symptoms of a vestibular schwannoma. These are differentiated using magnetic resonance angiography (MRA) or CT angiography.
    \item \textbf{Metastatic Disease and Inflammatory Lesions:} Metastases from lung, breast, or melanoma, or systemic inflammatory conditions like sarcoidosis, tuberculosis, or neurosyphilis can involve the CPA and present with cranial neuropathies. These are typically associated with rapid symptom progression and multifocal disease.
\end{itemize}

\section*{When to Seek Medical Care}
Patients experiencing any persistent or progressive unilateral hearing symptoms should undergo a formal medical evaluation. Early identification of a vestibular schwannoma is critical to preserving hearing and preventing long-term neurological deficits.

\subsection*{Routine Consultation}
A consultation with a general practitioner, audiologist, or an ear, nose, and throat (ENT) specialist should be scheduled if a patient experiences:
\begin{itemize}
    \item Progressive hearing loss that is noticeably worse in one ear.
    \item Persistent unilateral tinnitus (ringing, buzzing, or clicking in one ear).
    \item Unsteadiness, chronic imbalance, or mild dizziness that cannot be attributed to other clear causes.
    \item Mild, intermittent facial twitching or numbness.
\end{itemize}

\subsection*{Emergency Medical Attention}
Immediate emergency medical care must be sought if any of the following symptoms occur, as they may indicate acute hydrocephalus, brainstem compression, or intracranial haemorrhage. Patients or bystanders should immediately dial their national emergency service number, such as:
\begin{itemize}
    \item \textbf{911} in the United States and Canada,
    \item the local emergency number,
    \item \textbf{999} in the United Kingdom,
    \item \textbf{112} in the European Union and many other global regions.
\end{itemize}
Emergency services should be contacted for:
\begin{itemize}
    \item Sudden, severe, and progressive headache, particularly when accompanied by a stiff neck, nausea, and projectile vomiting.
    \item Acute onset of severe cerebellar ataxia, causing an inability to walk or stand without falling.
    \item Rapidly progressive weakness or paralysis of one side of the face.
    \item Sudden, unexplained changes in mental status, such as confusion, severe lethargy, obtundation, or loss of consciousness.
    \item Double vision (diplopia) or sudden loss of vision in one or both eyes.
    \item Difficulty swallowing (dysphagia) or speaking (dysarthria) that develops rapidly.
\end{itemize}

\section*{Diagnosis and Investigations}
The diagnostic pathway for vestibular schwannoma involves a combination of clinical assessment, audiometric testing, and advanced neuroimaging.
\begin{itemize}
    \item \textbf{Clinical Examination:} A comprehensive neurological examination is performed, with special emphasis on the cranial nerves. This includes assessing facial sensation (cranial nerve V), corneal reflex sensitivity, facial muscle symmetry and strength (cranial nerve VII), hearing and balance (cranial nerve VIII), and lower cranial nerve functions (cranial nerves IX, X, XI, and XII). Cerebellar testing, such as the Romberg test, tandem gait assessment, finger-to-nose testing, and rapid alternating movement evaluation, is essential to detect cerebellar compression.
    \item \textbf{Pure-Tone and Speech Audiometry:} Audiometry is the standard screening tool. The classic finding in vestibular schwannoma is an asymmetric sensorineural hearing loss, typically defined as a difference of 15 decibels (dB) or more at three contiguous frequencies between the ears. Additionally, speech discrimination scores are often disproportionately poor compared to the pure-tone average, representing retrocochlear pathology.
    \item \textbf{Auditory Brainstem Response (ABR):} ABR testing (or brainstem auditory evoked potentials) measures the electrical activity along the auditory pathway from the cochlea to the brainstem. While once a primary screening tool, ABR has been largely superseded by MRI due to a high false-negative rate for tumours smaller than 1 cm. However, it remains useful when MRI is contraindicated. Characteristic findings include a prolonged wave V latency or an increased I-V interpeak interval on the affected side.
    \item \textbf{Vestibular Testing:} Electronystagmography (ENG) or videonystagmography (VNG) may reveal a reduced caloric response (vestibular paresis) in the affected ear, confirming vestibular nerve dysfunction. Caloric testing demonstrates a unilateral weakness in approximately 80\% of patients.
    \item \textbf{Magnetic Resonance Imaging (MRI):} Contrast-enhanced MRI of the brain and internal auditory canals is the gold standard for diagnosing vestibular schwannoma. The imaging protocol must include thin-section T2-weighted sequences (such as FIESTA or CISS) and pre- and post-gadolinium contrast T1-weighted sequences. Vestibular schwannomas typically appear as intensely and homogeneously enhancing masses centred within the IAC, often with a spherical or ice-cream cone appearance extending into the CPA.
    \item \textbf{Computed Tomography (CT) Scan:} High-resolution CT of the temporal bones is indicated when MRI is contraindicated (e.g., in patients with non-compatible pacemakers, metallic implants, or severe claustrophobia). While CT is less sensitive for small intracanalicular tumours, it is excellent for demonstrating expansion and erosion of the bony internal auditory canal.
\end{itemize}

\section*{Management}
The management of vestibular schwannoma must be individualized based on tumour size, growth rate, patient age, baseline hearing status, general medical health, and patient preference. The three primary management strategies are active surveillance, radiation therapy, and surgical resection.
\begin{itemize}
    \item \textbf{Active Surveillance (``Watch and Wait''):} Because many vestibular schwannomas grow very slowly (average growth rate of 1 mm to 2 mm per year) or do not grow at all after detection, conservative management with serial imaging is highly appropriate for many patients. This strategy is preferred for small, intracanalicular tumours ($< 1.5$ cm), elderly patients with comorbidities, patients with poor baseline hearing who wish to delay intervention, or incidentally discovered asymptomatic lesions. Surveillance typically involves clinical evaluations and contrast-enhanced MRIs every 6 to 12 months initially, extending to every 1 to 2 years if stability is demonstrated. If documented growth occurs or new symptoms emerge, intervention is indicated.
    \item \textbf{Stereotactic Radiosurgery (SRS) and Radiotherapy:} Radiation therapy aims to arrest tumour growth by delivering highly targeted ionizing radiation, inducing cellular senescence and vascular sclerosis within the neoplasm. It does not remove the tumour but achieves long-term growth control in over 90\% of cases.
    \begin{itemize}
        \item \textit{Stereotactic Radiosurgery (SRS):} Typically delivered in a single fraction using a Gamma Knife, linear accelerator (LINAC), or CyberKnife system. SRS is highly effective for small- to medium-sized tumours ($< 3$ cm in maximum diameter) and is associated with low rates of facial nerve injury ($< 1\%$ to 2\%) and high rates of long-term hearing preservation.
        \item \textit{Fractionated Stereotactic Radiotherapy (FSRT):} Delivered over multiple daily sessions (typically 25 to 30 fractions). FSRT may be considered for slightly larger tumours or when there is a desire to maximize hearing preservation by exposing surrounding nerves to lower daily doses of radiation.
    \end{itemize}
    \item \textbf{Surgical Resection:} Microsurgical removal of the tumour is indicated for large, symptomatic tumours ($> 3$ cm) causing significant mass effect on the brainstem, tumours showing documented rapid growth, or patients with intractable symptoms. Surgery is performed under general anaesthesia with continuous intraoperative cranial nerve monitoring (facial and auditory nerves). There are three primary surgical approaches:
    \begin{itemize}
        \item \textit{Translabyrinthine Approach:} This approach involves drilling through the semicircular canals of the inner ear. It offers excellent, direct exposure of the internal auditory canal and CPA with minimal brainstem retraction. However, it results in complete, permanent hearing loss in the operated ear, making it suitable only for patients with poor baseline hearing.
        \item \textit{Retrosigmoid (Suboccipital) Approach:} This approach involves a craniotomy behind the ear and retraction of the cerebellum. It provides a wide view of the CPA and is suitable for tumours of any size. Importantly, it offers the potential to preserve hearing in patients with usable pre-operative hearing, although the rate of success depends on tumour size.
        \item \textit{Middle Cranial Fossa (MCF) Approach:} This approach involves a temporal craniotomy and accessing the IAC from above. It is technically demanding and reserved for small, purely intracanalicular tumours in patients with good hearing, with the primary goal of hearing preservation.
    \end{itemize}
\end{itemize}

\section*{Complications}
Complications can arise from the progressive growth of the untreated tumour or as a consequence of therapeutic interventions (surgery or radiation).
\begin{itemize}
    \item \textbf{Complications of Untreated Disease:}
    \begin{itemize}
        \item \textit{Irreversible Hearing Loss:} Complete sensorineural deafness in the affected ear.
        \item \textit{Obstructive Hydrocephalus:} Life-threatening accumulation of cerebrospinal fluid due to ventricular compression, requiring urgent ventriculoperitoneal shunt placement or tumour resection.
        \item \textit{Brainstem Compression:} Compression of vital cardiorespiratory centres, leading to severe neurological decline, coma, and death.
    \end{itemize}
    \item \textbf{Surgical Complications:}
    \begin{itemize}
        \item \textit{Facial Nerve Palsy:} Injury to the facial nerve during surgery can cause temporary or permanent facial paralysis, leading to cosmetic disfigurement, difficulty with speech and mastication, and incomplete eye closure (lagophthalmos), which carries a high risk of corneal ulceration and blindness.
        \item \textit{Cerebrospinal Fluid (CSF) Leak:} CSF can leak through the surgical wound or via the eustachian tube into the nose (rhinorrhea), occurring in 5\% to 10\% of surgeries. This requires wound revision, lumbar drain insertion, or surgical re-exploration to prevent meningitis.
        \item \textit{Vestibular Dysfunction:} Severe post-operative vertigo and persistent imbalance, which typically improves over weeks to months as the contralateral vestibular system compensates.
        \item \textit{Neurological Deficits:} Lower cranial nerve palsies, cerebellar injury, stroke, or meningitis.
    \end{itemize}
    \item \textbf{Radiation-Induced Complications:}
    \begin{itemize}
        \item \textit{Transient Tumour Swelling:} Radiation can cause inflammation and swelling within the first 6 to 18 months, which can temporarily worsen symptoms or cause hydrocephalus.
        \item \textit{Radiation-Induced Cranial Neuropathies:} Delayed facial weakness, facial numbness (trigeminal neuralgia), or progressive loss of remaining hearing.
        \item \textit{Secondary Malignancies:} A highly rare but potential risk of malignant transformation of the benign schwannoma or development of surrounding brain tumours decades after radiation.
    \end{itemize}
\end{itemize}

\section*{Prognosis and Outcomes}
The prognosis for patients with vestibular schwannoma is generally excellent regarding survival, given the benign nature of the tumour. The primary clinical challenges are related to the preservation of neurological function, particularly hearing, facial nerve function, and balance.

For patients undergoing active surveillance, approximately 30\% to 50\% will eventually require active treatment (surgery or radiation) within 5 years of diagnosis due to documented tumour growth or progressive symptoms. The remaining patients maintain stable tumour size with no significant change in neurological status.

Following stereotactic radiosurgery, the rate of long-term tumour control (defined as tumour shrinkage or stability) is between 90\% and 95\%. The rate of facial nerve preservation is excellent, exceeding 98\%. Long-term preservation of serviceable hearing is achieved in approximately 50\% to 70\% of patients at 5 years post-radiation, although this rate typically declines slowly over subsequent years.

Following microsurgery, complete tumour removal is achieved in the vast majority of cases, and recurrence is rare ($< 1\%$ to 2\% if complete resection is accomplished). The preservation of facial nerve function is closely linked to tumour size: for small tumours ($< 1.5$ cm), normal or near-normal facial function is preserved in over 90\% of cases; however, for large tumours ($> 3$ cm), the rate of long-term facial nerve function preservation drops to 50\% to 70\%. Hearing preservation is successful in approximately 30\% to 50\% of carefully selected candidates undergoing retrosigmoid or middle fossa resection. Post-operative balance recovery varies, but most patients achieve satisfactory compensation within 3 to 6 months with the help of specialized rehabilitation.

\section*{Prevention and Self-Care}
There are no established preventative measures to prevent the development of sporadic vestibular schwannomas, as their primary cause is genetic mutation. However, self-care, symptom management, and rehabilitation play a major role in maximizing quality of life.
\begin{itemize}
    \item \textbf{Vestibular Rehabilitation Therapy (VRT):} For patients experiencing chronic imbalance or recovering from surgery, VRT is highly beneficial. This specialized form of physical therapy uses exercises designed to promote brain compensation (habituation, gaze stabilization, and balance training) to adapt to the loss of vestibular input from one ear.
    \item \textbf{Corneal Protection:} Patients with facial nerve weakness or trigeminal nerve dysfunction often suffer from dry eyes and incomplete eye closure. Strict self-care is required to prevent corneal damage:
    \begin{itemize}
        \item Regular use of preservative-free artificial tears during the day.
        \item Application of thick ophthalmic ointment at bedtime.
        \item Taping the eyelid closed or wearing a moisture chamber patch overnight.
        \item Seeking prompt ophthalmology consultation if eye redness, pain, or visual changes occur.
    \end{itemize}
    \item \textbf{Hearing Conservation and Rehabilitation:} Protect the remaining hearing in the contralateral ear by avoiding loud noises and using ear protection. For the affected ear, rehabilitation options include:
    \begin{itemize}
        \item \textit{Conventional Hearing Aids:} Useful for mild-to-moderate hearing loss.
        \item \textit{Contralateral Routing of Signals (CROS) Hearing Aids:} Transmits sound from the deaf ear to the normal ear via a microphone and wireless transmitter.
        \item \textit{Bone-Anchored Hearing Systems (BAHS):} Uses bone conduction to transmit sound from the affected side through the skull to the functioning inner ear.
        \item \textit{Cochlear Implantation:} Increasingly considered, especially in patients with NF2 or bilateral hearing loss, to directly stimulate the auditory nerve.
    \end{itemize}
    \item \textbf{Support Groups and Psychological Self-Care:} Living with a chronic brain tumour, hearing loss, or facial disfigurement can cause significant emotional distress, anxiety, and depression. Joining patient support networks (such as the Acoustic Neuroma Association) and seeking professional counseling can improve coping mechanisms and mental health outcomes.
\end{itemize}

\section*{Sources and Further Reading}
\begin{itemize}
    \item \textbf{Acoustic Neuroma Association (ANA):} A leading patient advocacy group providing extensive educational resources, support networks, and clinical updates. Website: \url{https://www.anausa.org/}
    \item \textbf{National Institute on Deafness and Other Communication Disorders (NIDCD):} A division of the U.S. National Institutes of Health (NIH) providing authoritative fact sheets on vestibular schwannoma. Website: \url{https://www.nidcd.nih.gov/}
    \item \textbf{Mayo Clinic - Acoustic Neuroma:} A comprehensive clinical overview of symptoms, diagnosis, and treatment options from a world-renowned medical center. Website: \url{https://www.mayoclinic.org/diseases-conditions/acoustic-neuroma/symptoms-causes/syc-20354927}
    \item \textbf{National Health Service (NHS) UK - Acoustic Neuroma:} Public health information detailing symptoms, diagnostics, and treatment pathways in the United Kingdom. Website: \url{https://www.nhs.uk/conditions/acoustic-neuroma/}
    \item \textbf{The Lancet - Vestibular Schwannoma:} A high-impact peer-reviewed seminar article outlining the state-of-the-art management and scientific consensus on the disease. Website: \url{https://www.thelancet.com/}
    \item \textbf{Brain Tumour Research - Acoustic Neuroma:} A charity-driven resource offering information on brain tumour statistics, active research, and support. Website: \url{https://www.braintumourresearch.org/}
\end{itemize}
